% Handschrift: Seiten 14--18; Argument geometrisch statt Falltabelle (Autorwunsch).
\section{Polardarstellung und Darstellungswechsel}
\label{sec:complex_numbers_polar_form}
Statt durch zwei Koordinaten lässt sich eine komplexe Zahl durch ihren
Abstand vom Ursprung und ihre Richtung beschreiben. Diese Sicht bereitet
die geometrische Deutung von Produkten und Potenzen vor.

\begin{DefinitionBox}{Trigonometrische Polardarstellung}{complex_numbers_polar}
  Für $z\ne0$ heißen $r=|z|>0$ und ein Winkel $\varphi$ zur positiven
  reellen Achse \emph{Polarkoordinaten} \EnglishTerm{polar coordinates}.
  Der Winkel heißt \emph{Argument} \EnglishTerm{argument} oder Phasenwinkel.
  Die \emph{Polardarstellung} \EnglishTerm{polar form} lautet
  \[
    z=r(\cos\varphi+i\sin\varphi).
  \]
  Für $z=0$ ist kein Argument festgelegt.
  \begin{center}
    \CachedTikzFigure{content/resources/chapters/complex_numbers/tikz/complex_polar_coordinates.tex}
  \end{center}
\end{DefinitionBox}
Die rechtwinklige Zerlegung verbindet Polarkoordinaten und kartesische Koordinaten:
\[
  \cos\varphi=\frac{x}{r},\quad \sin\varphi=\frac{y}{r}
  \quad\Longrightarrow\quad
  x=r\cos\varphi,\quad y=r\sin\varphi.
\]
\begin{ExampleBox}{Betrag und Winkel in der Zahlenebene}{complex_numbers_polar_examples}
  Zeichnen Sie die Zahlen $z_1=3(\cos45^\circ+i\sin45^\circ)$ und
  $z_2=4(\cos150^\circ+i\sin150^\circ)$ ausgehend von Betrag und Winkel.
  \begin{center}
    \CachedTikzFigure{content/resources/chapters/complex_numbers/tikz/complex_polar_examples.tex}
  \end{center}
  Die Pfeillängen betragen $3$ beziehungsweise $4$; die Winkel werden
  von der positiven reellen Achse aus gemessen.
\end{ExampleBox}

\clearpage
\subsection{Volle Umdrehungen und besondere Winkel}
Eine volle Umdrehung verändert die Zahl nicht. Deshalb beschreiben
$\varphi$ und $\varphi+2k\pi$ mit $k\in\Z$ dieselbe Richtung:
\[
  \cos40^\circ+i\sin40^\circ
  =\cos400^\circ+i\sin400^\circ.
\]
Häufig wählt man einen Winkel in $[0,2\pi)$ beziehungsweise $[0^\circ,360^\circ)$.
Auch negative Winkel sind möglich; sie werden im Uhrzeigersinn gemessen.
\begin{ExampleBox}{Die vier Achsenrichtungen}{complex_numbers_trigonometric_axes}
  Lesen Sie die Werte aus der Lage auf dem Einheitskreis ab.
  \begin{center}
    \renewcommand{\arraystretch}{1.5}
    \begin{tabular}{@{}c@{\qquad}c@{\qquad}c@{}}
      $\varphi$ & $\cos\varphi+i\sin\varphi$ & Lage\\\midrule
      $0^\circ$ & $\Cloze{2cm}{1}$ & positive reelle Achse\\
      $90^\circ$ & $\Cloze{2cm}{i}$ & positive imaginäre Achse\\
      $180^\circ$ & $\Cloze{2cm}{-1}$ & negative reelle Achse\\
      $270^\circ$ & $\Cloze{2cm}{-i}$ & negative imaginäre Achse\\
      $360^\circ$ & $\Cloze{2cm}{1}$ & wieder die Ausgangsrichtung
    \end{tabular}
  \end{center}
\end{ExampleBox}

\subsection{Von der Polarform zur kartesischen Form}
Zur Umrechnung berechnen wir die beiden Komponenten $x=r\cos\varphi$
und $y=r\sin\varphi$ getrennt.
\begin{ExampleBox}{Komponenten aus Betrag und Winkel bestimmen}{complex_numbers_polar_to_cartesian}
  Bestimmen Sie Real- und Imaginärteil der Zahl
  $z=3(\cos45^\circ+i\sin45^\circ)$ aus dem vorigen Beispiel.
  \begin{ClozeArea}{5}
    \[
      x=3\cos45^\circ=\frac3{\sqrt2},\qquad
      y=3\sin45^\circ=\frac3{\sqrt2},\qquad
      z=\frac3{\sqrt2}(1+i).
    \]
  \end{ClozeArea}
\end{ExampleBox}

\clearpage
\subsection{Von der kartesischen Form zum Argument}
Der Betrag ist unmittelbar $r=\sqrt{x^2+y^2}$. Für den Winkel nutzen wir
bei $x\ne0$ die Beziehung
\[
  \tan\varphi=\frac{\sin\varphi}{\cos\varphi}
    =\frac{y/r}{x/r}=\frac{y}{x}.
\]
\begin{KeyPointBox}{Der Arkustangens kennt den Quadranten nicht}{complex_numbers_argument_geometry}
  $\arctan(y/x)$ liefert nur einen Winkel zwischen $-\pi/2$ und $\pi/2$.
  Gegenüberliegende Pfeile haben denselben Quotienten $y/x$ und werden
  dadurch nicht unterschieden.

  \textbf{Stellen Sie sich $z$ in der Zahlenebene vor oder zeichnen Sie es ein.}
  Prüfen Sie, ob der berechnete Winkel in die richtige Richtung zeigt.
  Zeigt er genau entgegengesetzt, addieren oder subtrahieren Sie
  $\pi$ beziehungsweise $180^\circ$ passend zur Skizze.

  Bei $x=0$ ist $y/x$ nicht definiert. Lesen Sie den Winkel direkt an der
  imaginären Achse ab: nach oben $\pi/2$, nach unten $-\pi/2$
  (gleichwertig $3\pi/2$). Für $z=0$ gibt es kein Argument.
\end{KeyPointBox}
\begin{ExampleBox}{Ein Punkt im dritten Quadranten}{complex_numbers_argument_example}
  Bestimmen Sie Betrag und einen passenden Winkel von $z=-1-i$.
  Vergleichen Sie die tatsächliche Richtung mit dem Ergebnis des Arkustangens.
  \par\smallskip
  \begin{minipage}[c]{0.47\linewidth}
    \CachedTikzFigure{content/resources/chapters/complex_numbers/tikz/complex_argument_quadrant.tex}
  \end{minipage}\hfill
  \begin{minipage}[c]{0.48\linewidth}
    Der Punkt liegt links unten. Dagegen zeigt
    $\arctan((-1)/(-1))=\pi/4$ nach rechts oben.
    Wir drehen diesen Pfeil um eine halbe Umdrehung.
    \begin{ClozeArea}[15]{6}
      \[
        r=\sqrt2,\qquad
        \varphi=\frac\pi4-\pi=-\frac{3\pi}{4}.
      \]
      Mit $+\pi$ ergibt sich gleichwertig $5\pi/4$.
    \end{ClozeArea}
  \end{minipage}
  \par\medskip
  Die trigonometrische Darstellung lautet damit
  \[
    z=\sqrt2\left(\cos\left(-\frac{3\pi}{4}\right)
                +i\sin\left(-\frac{3\pi}{4}\right)\right)
     =\sqrt2\left(\cos\frac{3\pi}{4}-i\sin\frac{3\pi}{4}\right).
  \]
\end{ExampleBox}
